%Example of Borsuk lattice diameter notion when non convex point set 

\begin{figure}[ht!]
    \centering
        \begin{tikzpicture}[thick, scale=0.8]\vspace{0.5cm}
        %\draw[line width=0.3mm, white](-3,-1.5) -- (6,-1.5);
    
        \foreach \x in {-1,...,5} {
            \foreach \y in {-1,...,1} {
                \node at (\x,\y) [circle, draw=lightgray, fill= lightgray, minimum width=2pt] {};   
            }
        }
    
        % \fill[line width=0.3mm, lightgray, opacity=0.6] (0,-1) -- (0,1) -- (4,0) -- cycle;
        \draw[line width=0.3mm, black] (0,-1) -- (0,1) -- (4,0) -- cycle;

         \foreach \x in {0,...,3} {
                \node at (\x,0) [circle, fill=darklav, minimum width=4pt] {};
            }
        \node at (0,-1) [circle, fill=orange, minimum width=4pt] {};
        \node at (0,1) [circle, fill=orange, minimum width=4pt] {};
        \node at (4,0) [circle, fill=orange, minimum width=4pt] {};

        \draw(2.5,0.75)node[black]{$P$};
        
        \draw(-0.5,0.2)node[darklav]{$S_2$};
        \draw(4.2,-0.5)node[orange]{$S_1$};
        

    \end{tikzpicture} %\vspace{-0.5cm}
    
    \caption{A Borsuk partition $P \cap \Z^d = S_1 \sqcup S_2$ where $\ldiam(S_1) < \ldiam(\conv(S_1))=\ldiam(P)$.}
    \label{fig:Borsuk-sets}
\end{figure}